% Part: many-valued-logic
% Chapter: infinite-valued-logics
% Section: introduction

\documentclass[../../../include/open-logic-section]{subfiles}

\begin{document}

\olfileid{mvl}{inf}{int}

\olsection{પરિચય}

શ્રેણિકમાં સત્યમૂલ્યોની સંખ્યા સાન્ત હોવી જરૂરી નથી. અનંત સત્યમૂલ્યોના
ગણ માટે સહજ પસંદગી $0$ અને~$1$ વચ્ચેની સંમેય સંખ્યાઓનો ગણ
$V_\infty = [0,1] \cap \Rat$ છે; એટલે કે,
\begin{align*}
    V_\infty & = \Setabs{\frac{n}{m}}{n,m \in \Nat \text{ અને } m>0 \text{ તથા } n\le m}.
\intertext{આ અનંત સત્યમૂલ્ય-ગણને ધ્યાનમાં લેતાં, તેના નીચેના ઉપગણો
પણ ધ્યાનમાં લેવા ઘણી વાર ઉપયોગી છે; અહીં $m\ge 2$ છે:}
V_m & = \Setabs{\frac{n}{m-1}}{n \in \Nat \text{ અને } n<m}
\intertext{દાખલા તરીકે, $V_5$ એ સમાન અંતરે આવેલાં $5$ સત્યમૂલ્યોનો ગણ છે,}
V_5 & = \{0, \frac{1}{4}, \frac{1}{2}, \frac{3}{4}, 1\}.
\end{align*}
\sourcecorrection{OLMV-011}{સ્થિર અંગ્રેજી મૂળના
અનંત સંમેય ગણમાં ભાગાકારના છેદ માટે શૂન્યને
બાકાત રાખતું નથી. શૂન્યનો ભાગાકાર અવ્યાખ્યાયિત
હોવાથી અહીં છેદ ધન હોય તે શરત સ્પષ્ટ કરી છે.}
\sourcecorrection{OLMV-012}{સ્થિર અંગ્રેજી મૂળના
સૂત્રમાં $V_m$ માટે અંશને $m$ સુધી મંજૂર કરે છે,
જેથી દર્શાવેલા $V_5$માં પાંચને બદલે છ મૂલ્યો અને
એક કરતાં મોટું મૂલ્ય મળે. અહીં અંશ $m$ કરતાં
નાનો રાખ્યો છે તથા છેદ શૂન્ય ન થાય તે માટે
$m\ge 2$ સ્પષ્ટ કર્યું છે.}
આ સત્યમૂલ્ય-ગણો પર આધારિત તર્કશાસ્ત્રોમાં સામાન્ય રીતે માત્ર
$1$ નિર્દિષ્ટ હોય છે; એટલે $V^+ = \{1\}$. બીજા શબ્દોમાં,
$1$ને (સંપૂર્ણ) સત્ય અને $0$ને સંપૂર્ણ અસત્યની ભૂમિકા આપીએ છીએ,
પરંતુ !!{formula}sનું મૂલ્ય~$V$માંનું કોઈ પણ મધ્યવર્તી મૂલ્ય હોઈ શકે.

આ ઉપરાંત $0$ અને~$1$ વચ્ચેની બધી \emph{વાસ્તવિક} સંખ્યાઓનો ગણ
$V_{[0,1]} = [0,1]$ અથવા $[0,1]$ના બીજા અનંત ઉપગણો પણ વિચારી શકાય.
આ સત્યમૂલ્ય-ગણ ધરાવતા તર્કશાસ્ત્રોને ઘણી વાર \emph{અસ્પષ્ટ}
તર્કશાસ્ત્રો કહે છે.

\end{document}
